% GENERATED FILE. Edit canonical lessons, then run npm run generate:books.
% canonical-source-hash: fnv1a64:daf84fab9f9c3d00
% canonical-lessons: PT-C54-oculos, PT-C54-caixa, PT-C54-mala, PT-C54-bilhete, PT-C54-mapa

\chapter{Glasses, Box, Suitcase, Ticket, Map}
\label{ch:pt-glasses-box-suitcase-ticket-map}
\hlchaptermodality{54}

\begin{chapteropening}
\textbf{By the end of this chapter:} \emph{I can name glasses, a box, a suitcase, a ticket and a map.}

Complete the last lesson of chapter 54: I can name glasses, a box, a suitcase, a ticket and a map.
\end{chapteropening}

\section[os óculos]{os óculos — glasses, spectacles}

\label{lesson:PT-C54-oculos}



\begin{quote}
Before the new one: say the Portuguese for a present, then the Portuguese for a watch.
\end{quote}

\subsection*{You'll want to know: os óculos}
\textbf{os óculos} — \textquotedblleft{}glasses, spectacles\textquotedblright{}. Plural in form.

\begin{cousinweb}
\textbf{óculos} $\leftarrow$ Latin \emph{oculus}, \textquotedblleft{}eye\textquotedblright{} (Lewis \& Short) — a learned borrowing from Latin. It is the plural of a learned word for a lens.
\end{cousinweb}

\begin{grammarlens}[title={What you've built}]
One more word for this chapter.
\end{grammarlens}

\subsection*{Your turn}
\begin{itemize}
\raggedright
  \item \emph{Say it:} \emph{os óculos}
  \item \emph{Say it:} \emph{os óculos}, once more
  \item \emph{Say it:} say \emph{os óculos}
  \item \emph{Recall:} say the Portuguese for a present, then the Portuguese for a watch, then say \emph{os óculos} again
\end{itemize}

\begin{tcolorbox}[breakable,colback=teal!4,colframe=teal!35!black,title={Before you move on}]
What does os óculos mean? (\textquotedblleft{}Glasses, spectacles\textquotedblright{}.) Say it once more.
\end{tcolorbox}
\section[a caixa]{a caixa — a box}

\label{lesson:PT-C54-caixa}



\begin{quote}
Before the new one: say the Portuguese for a watch, then the Portuguese for glasses.
\end{quote}

\subsection*{You'll want to know: a caixa}
\textbf{a caixa} — \textquotedblleft{}a box\textquotedblright{}.

\begin{cousinweb}
\textbf{caixa} — borrowed, most likely from Catalan \emph{caixa}, \textquotedblleft{}box\textquotedblright{}. EtymDB 2.1 traces it only to Old Portuguese \emph{cajxa}; it derives the Catalan word, through Old Provençal \emph{caissa}, from Latin \emph{capsa}, \textquotedblleft{}box\textquotedblright{}.
\end{cousinweb}

\begin{grammarlens}[title={What you've built}]
One more word for this chapter.
\end{grammarlens}

\subsection*{Your turn}
\begin{itemize}
\raggedright
  \item \emph{Say it:} \emph{a caixa}
  \item \emph{Say it:} \emph{a caixa}, once more
  \item \emph{Say it:} say \emph{a caixa}
  \item \emph{Recall:} say the Portuguese for a watch, then the Portuguese for glasses, then say \emph{a caixa} again
\end{itemize}

\begin{tcolorbox}[breakable,colback=teal!4,colframe=teal!35!black,title={Before you move on}]
What does a caixa mean? (\textquotedblleft{}A box\textquotedblright{}.) Say it once more.
\end{tcolorbox}
\section[a mala]{a mala — a suitcase}

\label{lesson:PT-C54-mala}



\begin{quote}
Before the new one: say the Portuguese for glasses, then the Portuguese for a box.
\end{quote}

\subsection*{You'll want to know: a mala}
\textbf{a mala} — \textquotedblleft{}a suitcase\textquotedblright{}.

\begin{cousinweb}
\textbf{mala} $\leftarrow$ French \emph{malle}, \textquotedblleft{}trunk\textquotedblright{} (EtymDB 2.1) — borrowed. The French word is of Germanic origin.
\end{cousinweb}

\begin{grammarlens}[title={What you've built}]
One more word for this chapter.
\end{grammarlens}

\subsection*{Your turn}
\begin{itemize}
\raggedright
  \item \emph{Say it:} \emph{a mala}
  \item \emph{Say it:} \emph{a mala}, once more
  \item \emph{Say it:} say \emph{a mala}
  \item \emph{Recall:} say the Portuguese for glasses, then the Portuguese for a box, then say \emph{a mala} again
\end{itemize}

\begin{tcolorbox}[breakable,colback=teal!4,colframe=teal!35!black,title={Before you move on}]
What does a mala mean? (\textquotedblleft{}A suitcase\textquotedblright{}.) Say it once more.
\end{tcolorbox}
\section[o bilhete]{o bilhete — a ticket}

\label{lesson:PT-C54-bilhete}



\begin{quote}
Before the new one: say the Portuguese for a box, then the Portuguese for a suitcase.
\end{quote}

\subsection*{You'll want to know: o bilhete}
\textbf{o bilhete} — \textquotedblleft{}a ticket\textquotedblright{}.

\begin{cousinweb}
\textbf{bilhete} $\leftarrow$ French \emph{billet}, \textquotedblleft{}note\textquotedblright{} (EtymDB 2.1) — borrowed. Behind the French word is Latin \emph{bulla}, a seal on a document.
\end{cousinweb}

\begin{grammarlens}[title={What you've built}]
One more word for this chapter.
\end{grammarlens}

\subsection*{Your turn}
\begin{itemize}
\raggedright
  \item \emph{Say it:} \emph{o bilhete}
  \item \emph{Say it:} \emph{o bilhete}, once more
  \item \emph{Say it:} say \emph{o bilhete}
  \item \emph{Recall:} say the Portuguese for a box, then the Portuguese for a suitcase, then say \emph{o bilhete} again
\end{itemize}

\begin{tcolorbox}[breakable,colback=teal!4,colframe=teal!35!black,title={Before you move on}]
What does o bilhete mean? (\textquotedblleft{}A ticket\textquotedblright{}.) Say it once more.
\end{tcolorbox}
\section[o mapa]{o mapa — a map}

\label{lesson:PT-C54-mapa}



\begin{quote}
Before the new one: say the Portuguese for a suitcase, then the Portuguese for a ticket.
\end{quote}

\subsection*{You'll want to know: o mapa}
\textbf{o mapa} — \textquotedblleft{}a map\textquotedblright{}. Masculine, despite the -a.

\begin{cousinweb}
\textbf{mapa} $\leftarrow$ Latin \emph{mappa}, \textquotedblleft{}napkin, cloth\textquotedblright{} (EtymDB 2.1) — a learned borrowing from Latin. A map was a \emph{mappa mundi}, a \textquotedblleft{}cloth of the world\textquotedblright{}.
\end{cousinweb}

\begin{grammarlens}[title={What you've built}]
One more word for this chapter.
\end{grammarlens}

\subsection*{Your turn}
\begin{itemize}
\raggedright
  \item \emph{Say it:} \emph{o mapa}
  \item \emph{Say it:} \emph{o mapa}, once more
  \item \emph{Say it:} say \emph{o mapa}
  \item \emph{Recall:} say the Portuguese for a suitcase, then the Portuguese for a ticket, then say \emph{o mapa} again
\end{itemize}

\begin{tcolorbox}[breakable,colback=teal!4,colframe=teal!35!black,title={Before you move on}]
What does o mapa mean? (\textquotedblleft{}A map\textquotedblright{}.) Say it once more.
\end{tcolorbox}
